\uuid{r9Di}
\titre{Résolution d'équation différentielle du premier ordre}
\auteur{bodin ; croizat ; sacre}
\organisation{exo7}
\licence{CC BY-NC-SA 4.0}
\datecreate{2022-03-08}
\niveau{L1}
\module{Analyse}
\chapitre{Equation différentielle}
\sousChapitre{Résolution d'équation différentielle du premier ordre}
\difficulte{3}
\theme{}
\exo7qcmid{hilisit.equations-differentielles.037}

\contenu{
\begin{question}
\qclassification{225.01}
Soit la solution de l'équation différentielle $y'= 2y$ vérifiant $y(0) = -1$.
Quelles sont les affirmations vraies ?
\begin{answers}
  \good{La solution est toujours négative.}
  \good{La solution est une fonction décroissante.}
  \bad{La pente de la tangente en $x=0$ vaut $1$.}
  \good{La pente de la tangente en $x=1$ vaut $-2e^2$.}
\end{answers}
\begin{explanations}
Les solutions générales sont $y(x) = Ce^{2x}$. Comme $y(0)=-1$ alors $C = -1$.
La solution est donc  $f(x) = -e^{2x}$.
La pente de la tangente en $x_0$ est donnée par $f'(x_0)$.
Comme $f(0)=-1$ alors $f'(0) = -2$, la pente de la tangente en $x=0$ vaut $-2$.
De façon générale, comme $f(x) = -e^{2x}$, alors $f'(x) = -2e^{2x}$ qui est une fonction toujours négative : ainsi $f$ est une fonction décroissante. La pente de sa tangente en $x=1$ vaut bien $f'(1) = -2e^2$.
\end{explanations}
\end{question}
}
